% GENERATED FILE. Edit canonical lessons, then run npm run generate:books.
% canonical-source-hash: fnv1a64:99b6d2a329a50b4c
% canonical-lessons: RU-C370-uslovie, RU-C370-pri-uslovii, RU-C370-v-sluchae, RU-C370-v-protivnom-sluchae, RU-C370-dazhe-esli

\chapter{On one condition}
\label{ch:ru-on-one-condition}
\hlchaptermodality{370}

\begin{chapteropening}
\textbf{By the end of this chapter:} \emph{I can set a condition: on condition that, in case of, otherwise and even if, and know what each one takes.}

Complete the last lesson of chapter 370: I can set a condition: on condition that, in case of, otherwise and even if, and know what each one takes.
\end{chapteropening}

\section[\texorpdfstring{\ru{условие} (uslóviye)}{условие (uslóviye)}]{\ru{условие} (uslóviye) — a condition}

\label{lesson:RU-C370-uslovie}



\begin{quote}
Before the new one: say the Russian for to recycle, then the Russian for to look after.
\end{quote}

\subsection*{You'll want to know: \ru{условие}}
\textbf{\ru{условие}} (\emph{uslóviye}) — \textquotedblleft{}a condition\textquotedblright{}.

\begin{quote}
\ru{У} \ru{меня} \ru{одно} \ru{условие}: \ru{мы} \ru{вернёмся} \ru{до} \ru{десяти}.
\end{quote}

\begin{quote}
\emph{U menyá odnó uslóviye: my vernyómsya do desyatí.}
\end{quote}

\begin{quote}
I have one condition: we come back before ten.
\end{quote}

Neuter. The plural \emph{\ru{условия}} also means terms or circumstances: \emph{\ru{условия} \ru{работы}}.

\begin{cousinweb}
\textbf{\ru{условие}} is built on the old verb \emph{\ru{условиться}} \textquotedblleft{}to agree on terms\textquotedblright{}, from \emph{\ru{у}-} and \emph{\ru{слово}} \textquotedblleft{}word\textquotedblright{}, inherited from Proto-Slavic \emph{*slovo} (EtymDB). A condition is something settled in words.
\end{cousinweb}

\begin{grammarlens}[title={What you've built}]
One more word for telling, explaining and describing at B1.
\end{grammarlens}

\subsection*{Your turn}
\begin{itemize}
\raggedright
  \item \emph{Say it:} \emph{uslóviye}
  \item \emph{Say it:} \emph{uslóviye}, once more
  \item \emph{Say it:} a sentence of your own with \emph{uslóviye}
  \item \emph{Recall:} say the Russian for to recycle, then the Russian for to look after, then say \emph{uslóviye} again
\end{itemize}

\begin{tcolorbox}[breakable,colback=teal!4,colframe=teal!35!black,title={Before you move on}]
What does \ru{условие} mean? (\textquotedblleft{}A condition\textquotedblright{}.) Say it once more.
\end{tcolorbox}
\section[\texorpdfstring{\ru{при} \ru{условии} (pri uslóvii)}{при условии (pri uslóvii)}]{\ru{при} \ru{условии} (pri uslóvii) — on condition, provided (that)}

\label{lesson:RU-C370-pri-uslovii}



\begin{quote}
Before the new one: say the Russian for to look after, then the Russian for a condition.
\end{quote}

\subsection*{You'll want to know: \ru{при} \ru{условии}}
\textbf{\ru{при} \ru{условии}} (\emph{pri uslóvii}) — \textquotedblleft{}on condition, provided (that)\textquotedblright{}.

\begin{quote}
\ru{Я} \ru{пойду} \ru{при} \ru{условии}, \ru{что} \ru{ты} \ru{пойдёшь} \ru{со} \ru{мной}.
\end{quote}

\begin{quote}
\emph{Ya poydú pri uslóvii, chto ty poydyósh so mnoy.}
\end{quote}

\begin{quote}
I will go on condition that you come with me.
\end{quote}

\emph{\ru{При} \ru{условии}, \ru{что}…} takes an ordinary indicative clause. With a noun it takes the genitive: \emph{\ru{при} \ru{условии} \ru{оплаты}}.

\begin{cousinweb}
\textbf{\ru{при} \ru{условии}} is \emph{\ru{при}} \textquotedblleft{}at, under\textquotedblright{} with the prepositional of \emph{\ru{условие}} \textquotedblleft{}a condition\textquotedblright{}, which is built on \emph{\ru{слово}} \textquotedblleft{}word\textquotedblright{}, Proto-Slavic \emph{*slovo} (EtymDB).
\end{cousinweb}

\begin{grammarlens}[title={What you've built}]
One more word for telling, explaining and describing at B1.
\end{grammarlens}

\subsection*{Your turn}
\begin{itemize}
\raggedright
  \item \emph{Say it:} \emph{pri uslóvii}
  \item \emph{Say it:} \emph{pri uslóvii}, once more
  \item \emph{Say it:} a sentence of your own with \emph{pri uslóvii}
  \item \emph{Recall:} say the Russian for to look after, then the Russian for a condition, then say \emph{pri uslóvii} again
\end{itemize}

\begin{tcolorbox}[breakable,colback=teal!4,colframe=teal!35!black,title={Before you move on}]
What does \ru{при} \ru{условии} mean? (\textquotedblleft{}On condition, provided (that)\textquotedblright{}.) Say it once more.
\end{tcolorbox}
\section[\texorpdfstring{\ru{в} \ru{случае} (v slúchaye)}{в случае (v slúchaye)}]{\ru{в} \ru{случае} (v slúchaye) — in case of, in the event of}

\label{lesson:RU-C370-v-sluchae}



\begin{quote}
Before the new one: say the Russian for a condition, then the Russian for on condition.
\end{quote}

\subsection*{You'll want to know: \ru{в} \ru{случае}}
\textbf{\ru{в} \ru{случае}} (\emph{v slúchaye}) — \textquotedblleft{}in case of, in the event of\textquotedblright{}.

\begin{quote}
\ru{В} \ru{случае} \ru{дождя} \ru{концерт} \ru{будет} \ru{в} \ru{зале}.
\end{quote}

\begin{quote}
\emph{V slúchaye dozhdyá kontsért búdet v zále.}
\end{quote}

\begin{quote}
In case of rain the concert will be in the hall.
\end{quote}

\emph{\ru{В} \ru{случае}} + genitive: \emph{\ru{в} \ru{случае} \ru{дождя}}. Before a clause, \emph{\ru{в} \ru{случае} \ru{если}…}

\begin{cousinweb}
\textbf{\ru{в} \ru{случае}} is \emph{\ru{в}} \textquotedblleft{}in\textquotedblright{} with the prepositional of \emph{\ru{случай}} \textquotedblleft{}an event, a chance\textquotedblright{}, inherited from Proto-Slavic \emph{*sŭlučajĭ}, from \emph{*lučiti}, tentatively \textquotedblleft{}to find, to reach\textquotedblright{} (EtymDB).
\end{cousinweb}

\begin{grammarlens}[title={What you've built}]
One more word for telling, explaining and describing at B1.
\end{grammarlens}

\subsection*{Your turn}
\begin{itemize}
\raggedright
  \item \emph{Say it:} \emph{v slúchaye}
  \item \emph{Say it:} \emph{v slúchaye}, once more
  \item \emph{Say it:} a sentence of your own with \emph{v slúchaye}
  \item \emph{Recall:} say the Russian for a condition, then the Russian for on condition, then say \emph{v slúchaye} again
\end{itemize}

\begin{tcolorbox}[breakable,colback=teal!4,colframe=teal!35!black,title={Before you move on}]
What does \ru{в} \ru{случае} mean? (\textquotedblleft{}In case of, in the event of\textquotedblright{}.) Say it once more.
\end{tcolorbox}
\section[\texorpdfstring{\ru{в} \ru{противном} \ru{случае} (v protívnom …)}{в противном случае (v protívnom …)}]{\ru{в} \ru{противном} \ru{случае} (v protívnom slúchaye) — otherwise, if not}

\label{lesson:RU-C370-v-protivnom-sluchae}



\begin{quote}
Before the new one: say the Russian for on condition, then the Russian for in case of.
\end{quote}

\subsection*{You'll want to know: \ru{в} \ru{противном} \ru{случае}}
\textbf{\ru{в} \ru{противном} \ru{случае}} (\emph{v protívnom slúchaye}) — \textquotedblleft{}otherwise, if not\textquotedblright{}.

\begin{quote}
\ru{Нужно} \ru{выйти} \ru{сейчас}, \ru{в} \ru{противном} \ru{случае} \ru{мы} \ru{опоздаем}.
\end{quote}

\begin{quote}
\emph{Núzhno výyti seychás, v protívnom slúchaye my opozdáyem.}
\end{quote}

\begin{quote}
We need to leave now, otherwise we will be late.
\end{quote}

Formal; in conversation people say \emph{\ru{иначе}} or \emph{\ru{а} \ru{то}}.

\begin{cousinweb}
\textbf{\ru{в} \ru{противном} \ru{случае}} is \textquotedblleft{}in the opposite case\textquotedblright{}: \emph{\ru{противный}} \textquotedblleft{}opposite\textquotedblright{}, built on \emph{\ru{против}} \textquotedblleft{}against\textquotedblright{}, and \emph{\ru{случай}} \textquotedblleft{}a case\textquotedblright{}, from Proto-Slavic \emph{*sŭlučajĭ} (EtymDB). Today \emph{\ru{противный}} on its own usually means nasty.
\end{cousinweb}

\begin{grammarlens}[title={What you've built}]
One more word for telling, explaining and describing at B1.
\end{grammarlens}

\subsection*{Your turn}
\begin{itemize}
\raggedright
  \item \emph{Say it:} \emph{v protívnom slúchaye}
  \item \emph{Say it:} \emph{v protívnom slúchaye}, once more
  \item \emph{Say it:} a sentence of your own with \emph{v protívnom slúchaye}
  \item \emph{Recall:} say the Russian for on condition, then the Russian for in case of, then say \emph{v protívnom slúchaye} again
\end{itemize}

\begin{tcolorbox}[breakable,colback=teal!4,colframe=teal!35!black,title={Before you move on}]
What does \ru{в} \ru{противном} \ru{случае} mean? (\textquotedblleft{}Otherwise, if not\textquotedblright{}.) Say it once more.
\end{tcolorbox}
\section[\texorpdfstring{\ru{даже} \ru{если} (dázhe yésli)}{даже если (dázhe yésli)}]{\ru{даже} \ru{если} (dázhe yésli) — even if}

\label{lesson:RU-C370-dazhe-esli}



\begin{quote}
Before the new one: say the Russian for in case of, then the Russian for otherwise.
\end{quote}

\subsection*{You'll want to know: \ru{даже} \ru{если}}
\textbf{\ru{даже} \ru{если}} (\emph{dázhe yésli}) — \textquotedblleft{}even if\textquotedblright{}.

\begin{quote}
\ru{Я} \ru{приду}, \ru{даже} \ru{если} \ru{будет} \ru{дождь}.
\end{quote}

\begin{quote}
\emph{Ya pridú, dázhe yésli búdet dozhd.}
\end{quote}

\begin{quote}
I will come even if it rains.
\end{quote}

Like \emph{\ru{если}}, it takes an ordinary clause in any tense: \emph{\ru{даже} \ru{если} \ru{он} \ru{позвонит}}. With \emph{\ru{бы}} it turns hypothetical: \emph{\ru{даже} \ru{если} \ru{бы} \ru{он} \ru{позвонил}}.

\begin{cousinweb}
\textbf{\ru{даже} \ru{если}} is built from \emph{\ru{даже}} \textquotedblleft{}even\textquotedblright{}, made of \emph{\ru{да}} and the particle \emph{\ru{же}}, and \emph{\ru{если}} \textquotedblleft{}if\textquotedblright{}, which is usually explained as \emph{\ru{есть}} \textquotedblleft{}is\textquotedblright{} with the particle \emph{\ru{ли}}.
\end{cousinweb}

\begin{grammarlens}[title={What you've built}]
One more word for telling, explaining and describing at B1.
\end{grammarlens}

\subsection*{Your turn}
\begin{itemize}
\raggedright
  \item \emph{Say it:} \emph{dázhe yésli}
  \item \emph{Say it:} \emph{dázhe yésli}, once more
  \item \emph{Say it:} a sentence of your own with \emph{dázhe yésli}
  \item \emph{Recall:} say the Russian for in case of, then the Russian for otherwise, then say \emph{dázhe yésli} again
\end{itemize}

\begin{tcolorbox}[breakable,colback=teal!4,colframe=teal!35!black,title={Before you move on}]
What does \ru{даже} \ru{если} mean? (\textquotedblleft{}Even if\textquotedblright{}.) Say it once more.
\end{tcolorbox}
